% GENERATED FILE. Edit canonical lessons, then run npm run generate:books.
% canonical-source-hash: fnv1a64:a2a0eb616a6f7d4e
% canonical-lessons: TA-C494-pankali, TA-C494-paramparai, TA-C494-talaimurai, TA-C494-cantati, TA-C494-varicu

\chapter{Kin and Lineage}
\label{ch:ta-kin-and-lineage}
\hlchaptermodality{494}

\begin{chapteropening}
\textbf{By the end of this chapter:} \emph{I can talk about kinsmen, lineage, generations, descendants and heirs.}

Complete the last lesson of chapter 494: I can talk about kinsmen, lineage, generations, descendants and heirs.
\end{chapteropening}

\section[\texorpdfstring{\ta{பங்காளி} (paṅkāḷi)}{பங்காளி (paṅkāḷi)}]{\ta{பங்காளி} (paṅkāḷi) — a paternal kinsman; a partner}

\label{lesson:TA-C494-pankali}



\begin{quote}
Before the new one: say the Tamil for a co-parent-in-law, then the Tamil for a mother’s younger sister.
\end{quote}

\subsection*{You'll want to know: \ta{பங்காளி}}
\textbf{\ta{பங்காளி}} — \emph{paṅkāḷi} — \textquotedblleft{}a paternal kinsman; a partner\textquotedblright{}.

\begin{quote}
\ta{திருவிழாவுக்கு} \ta{எல்லாப்} \ta{பங்காளிகளும்} \ta{வந்தார்கள்}.
\end{quote}

\begin{quote}
\emph{tiruviḻāvukku ellāp paṅkāḷikaḷum vantārkaḷ.}
\end{quote}

\begin{quote}
All our kinsmen came to the festival.
\end{quote}

Plural: \ta{பங்காளிகள்}, \emph{paṅkāḷikaḷ}.

\ta{பங்கு} \textquotedblleft{}a share\textquotedblright{} + -\ta{ஆளி} \textquotedblleft{}one who\textquotedblright{}: those who share an inheritance. Also a business partner.

\begin{cousinweb}
\textbf{\ta{பங்காளி}} (\emph{paṅkāḷi}) is a native compound of \ta{பங்கு} \emph{paṅku}, \textquotedblleft{}a share\textquotedblright{} (DEDR 3808), and -\ta{ஆளி}, \textquotedblleft{}one who has\textquotedblright{}.
\end{cousinweb}

\begin{grammarlens}[title={What you've built}]
One more word for describing what you lived and felt.
\end{grammarlens}

\subsection*{Your turn}
\begin{itemize}
\raggedright
  \item \emph{Say it:} \emph{paṅkāḷi}
  \item \emph{Say it:} \emph{paṅkāḷi}, once more
  \item \emph{Say it:} \emph{tiruviḻāvukku ellāp paṅkāḷikaḷum vantārkaḷ.} — All our kinsmen came to the festival.
  \item \emph{Recall:} say the Tamil for a co-parent-in-law, then the Tamil for a mother’s younger sister, then say \emph{paṅkāḷi} again
\end{itemize}

\begin{tcolorbox}[breakable,colback=teal!4,colframe=teal!35!black,title={Before you move on}]
What does \ta{பங்காளி} mean? (\textquotedblleft{}A paternal kinsman; a partner\textquotedblright{}.) Say it once more.
\end{tcolorbox}
\section[\texorpdfstring{\ta{பரம்பரை} (paramparai)}{பரம்பரை (paramparai)}]{\ta{பரம்பரை} (paramparai) — lineage; heredity}

\label{lesson:TA-C494-paramparai}



\begin{quote}
Before the new one: say the Tamil for a mother’s younger sister, then the Tamil for a paternal kinsman.
\end{quote}

\subsection*{You'll want to know: \ta{பரம்பரை}}
\textbf{\ta{பரம்பரை}} — \emph{paramparai} — \textquotedblleft{}lineage; heredity\textquotedblright{}.

\begin{quote}
\ta{சர்க்கரை} \ta{நோய்} \ta{எங்கள்} \ta{குடும்பத்தில்} \ta{பரம்பரையாக} \ta{வருகிறது}.
\end{quote}

\begin{quote}
\emph{carkkarai nōy eṅkaḷ kuṭumpattil paramparaiyāka varukiṟatu.}
\end{quote}

\begin{quote}
Diabetes runs in our family.
\end{quote}

\ta{பரம்பரையாக}, from generation to generation.

\begin{cousinweb}
\textbf{\ta{பரம்பரை}} (\emph{paramparai}) is borrowed from Sanskrit \emph{paramparā}, \textquotedblleft{}an uninterrupted series, succession\textquotedblright{} (Monier-Williams).
\end{cousinweb}

\begin{grammarlens}[title={What you've built}]
One more word for describing what you lived and felt.
\end{grammarlens}

\subsection*{Your turn}
\begin{itemize}
\raggedright
  \item \emph{Say it:} \emph{paramparai}
  \item \emph{Say it:} \emph{paramparai}, once more
  \item \emph{Say it:} \emph{carkkarai nōy eṅkaḷ kuṭumpattil paramparaiyāka varukiṟatu.} — Diabetes runs in our family.
  \item \emph{Recall:} say the Tamil for a mother’s younger sister, then the Tamil for a paternal kinsman, then say \emph{paramparai} again
\end{itemize}

\begin{tcolorbox}[breakable,colback=teal!4,colframe=teal!35!black,title={Before you move on}]
What does \ta{பரம்பரை} mean? (\textquotedblleft{}Lineage; heredity\textquotedblright{}.) Say it once more.
\end{tcolorbox}
\section[\texorpdfstring{\ta{தலைமுறை} (talaimuṟai)}{தலைமுறை (talaimuṟai)}]{\ta{தலைமுறை} (talaimuṟai) — a generation}

\label{lesson:TA-C494-talaimurai}



\begin{quote}
Before the new one: say the Tamil for a paternal kinsman, then the Tamil for lineage.
\end{quote}

\subsection*{You'll want to know: \ta{தலைமுறை}}
\textbf{\ta{தலைமுறை}} — \emph{talaimuṟai} — \textquotedblleft{}a generation\textquotedblright{}.

\begin{quote}
\ta{எங்கள்} \ta{வீட்டில்} \ta{மூன்று} \ta{தலைமுறைகள்} \ta{ஒன்றாக} \ta{வாழ்கின்றன}.
\end{quote}

\begin{quote}
\emph{eṅkaḷ vīṭṭil mūṉṟu talaimuṟaikaḷ oṉṟāka vāḻkiṉṟaṉa.}
\end{quote}

\begin{quote}
Three generations live together in our house.
\end{quote}

Plural: \ta{தலைமுறைகள்}, \emph{talaimuṟaikaḷ}.

\ta{தலை} \textquotedblleft{}head\textquotedblright{} + \ta{முறை} \textquotedblleft{}turn, order\textquotedblright{}.

\begin{cousinweb}
\textbf{\ta{தலைமுறை}} (\emph{talaimuṟai}) is a native compound of \ta{தலை} \emph{talai}, \textquotedblleft{}head\textquotedblright{} (DEDR 3103), and \ta{முறை} \emph{muṟai}, \textquotedblleft{}order, turn\textquotedblright{} (DEDR 5015).
\end{cousinweb}

\begin{grammarlens}[title={What you've built}]
One more word for describing what you lived and felt.
\end{grammarlens}

\subsection*{Your turn}
\begin{itemize}
\raggedright
  \item \emph{Say it:} \emph{talaimuṟai}
  \item \emph{Say it:} \emph{talaimuṟai}, once more
  \item \emph{Say it:} \emph{eṅkaḷ vīṭṭil mūṉṟu talaimuṟaikaḷ oṉṟāka vāḻkiṉṟaṉa.} — Three generations live together in our house.
  \item \emph{Recall:} say the Tamil for a paternal kinsman, then the Tamil for lineage, then say \emph{talaimuṟai} again
\end{itemize}

\begin{tcolorbox}[breakable,colback=teal!4,colframe=teal!35!black,title={Before you move on}]
What does \ta{தலைமுறை} mean? (\textquotedblleft{}A generation\textquotedblright{}.) Say it once more.
\end{tcolorbox}
\section[\texorpdfstring{\ta{சந்ததி} (cantati)}{சந்ததி (cantati)}]{\ta{சந்ததி} (cantati) — offspring; descendants}

\label{lesson:TA-C494-cantati}



\begin{quote}
Before the new one: say the Tamil for lineage, then the Tamil for a generation.
\end{quote}

\subsection*{You'll want to know: \ta{சந்ததி}}
\textbf{\ta{சந்ததி}} — \emph{cantati} — \textquotedblleft{}offspring; descendants\textquotedblright{}.

\begin{quote}
\ta{நம்} \ta{சந்ததிக்காக} \ta{இயற்கையைக்} \ta{காக்க} \ta{வேண்டும்}.
\end{quote}

\begin{quote}
\emph{nam cantatikkāka iyaṟkaiyaik kākka vēṇṭum.}
\end{quote}

\begin{quote}
We must protect nature for our descendants.
\end{quote}

A collective noun for all who come after.

\begin{cousinweb}
\textbf{\ta{சந்ததி}} (\emph{cantati}) is borrowed from Sanskrit \emph{saṃtati}, \textquotedblleft{}continuity, a line of descent\textquotedblright{} (Monier-Williams).
\end{cousinweb}

\begin{grammarlens}[title={What you've built}]
One more word for describing what you lived and felt.
\end{grammarlens}

\subsection*{Your turn}
\begin{itemize}
\raggedright
  \item \emph{Say it:} \emph{cantati}
  \item \emph{Say it:} \emph{cantati}, once more
  \item \emph{Say it:} \emph{nam cantatikkāka iyaṟkaiyaik kākka vēṇṭum.} — We must protect nature for our descendants.
  \item \emph{Recall:} say the Tamil for lineage, then the Tamil for a generation, then say \emph{cantati} again
\end{itemize}

\begin{tcolorbox}[breakable,colback=teal!4,colframe=teal!35!black,title={Before you move on}]
What does \ta{சந்ததி} mean? (\textquotedblleft{}Offspring; descendants\textquotedblright{}.) Say it once more.
\end{tcolorbox}
\section[\texorpdfstring{\ta{வாரிசு} (vāricu)}{வாரிசு (vāricu)}]{\ta{வாரிசு} (vāricu) — an heir}

\label{lesson:TA-C494-varicu}



\begin{quote}
Before the new one: say the Tamil for a generation, then the Tamil for offspring.
\end{quote}

\subsection*{You'll want to know: \ta{வாரிசு}}
\textbf{\ta{வாரிசு}} — \emph{vāricu} — \textquotedblleft{}an heir\textquotedblright{}.

\begin{quote}
\ta{அவருடைய} \ta{மகன்தான்} \ta{அந்தக்} \ta{கடைக்கு} \ta{வாரிசு}.
\end{quote}

\begin{quote}
\emph{avaruṭaiya makaṉtāṉ antak kaṭaikku vāricu.}
\end{quote}

\begin{quote}
His son is the heir to that shop.
\end{quote}

The thing inherited takes -\ta{க்கு}.

\begin{cousinweb}
\textbf{\ta{வாரிசு}} (\emph{vāricu}) is usually traced to Arabic \emph{wāriṯ}, \textquotedblleft{}an heir\textquotedblright{}, reaching Tamil through Hindustani; our sources do not document the path.
\end{cousinweb}

\begin{grammarlens}[title={What you've built}]
One more word for describing what you lived and felt.
\end{grammarlens}

\subsection*{Your turn}
\begin{itemize}
\raggedright
  \item \emph{Say it:} \emph{vāricu}
  \item \emph{Say it:} \emph{vāricu}, once more
  \item \emph{Say it:} \emph{avaruṭaiya makaṉtāṉ antak kaṭaikku vāricu.} — His son is the heir to that shop.
  \item \emph{Recall:} say the Tamil for a generation, then the Tamil for offspring, then say \emph{vāricu} again
\end{itemize}

\begin{tcolorbox}[breakable,colback=teal!4,colframe=teal!35!black,title={Before you move on}]
What does \ta{வாரிசு} mean? (\textquotedblleft{}An heir\textquotedblright{}.) Say it once more.
\end{tcolorbox}
